\section{Historical Fixed-Vector Comparisons}
\label{app:historical-vectors}
This secondary analysis compares a selected within-math archetype (A2), matched SEAL and random controls, and two saved biases from historical reward-learning runs, initialized from zero or A2. Only the layer-16 bias was optimized, using 301 math items disjoint from the 199 evaluation items, with one training seed per initialization. Evaluation uses corrected grades. The source audit identifies both an earlier reward parser and a temperature-variable defect; no per-run source hash or rollout log establishes their exact effect on training. We therefore analyze these as fixed interventions, not evidence for a reproducible tendency of correctly implemented RL (Appendix~\ref{app:audit}).

\begin{figure*}[t]
\centering\includegraphics[width=.95\textwidth,height=.47\textheight,keepaspectratio]{figures/fig05_steering_and_rl.pdf}
\caption{\textbf{Interventions differ in behavioral magnitude and performance despite similar shift directions.} Qwen3.5-9B, including saved historical RL vectors. (A) Paired first-8,192-token shifts: Q, questioning; P, perspective shift; C, conflict; R, reconciliation; V, verification; B, backtracking; S, subgoals; K, backward chaining. (B) Similarity to A2 in residual and mean behavior-change space; behavioral whiskers jointly bootstrap prompts. (C) Natural-completion accuracy, including matched SEAL and random controls. (D) Historical forced-answer differences with 95\% task-paired intervals; random is included, SEAL was not evaluated under this protocol. The RL training audit and answer-elicitation compute limit interpretation (Appendix~\ref{app:audit}).}
\label{fig:rl}
\end{figure*}
The saved zero-start RL vector has residual cosine approximately $-0.02$ with A2, but cosine 0.938 [0.839, 0.974] between their mean eight-rate shifts. The RL shift has only 0.412 times A2's magnitude. Both increase subgoals and reduce question--answer turns, conflict, and reconciliation. The matched SEAL direction is similarly aligned with A2 (0.934 [0.877, 0.968]), so the resemblance is not distinctive to reward learning. At a 1,024-token window, the A2--RL cosine is 0.740 [0.231, 0.911]. Joint prompt bootstraps, shorter windows, and disjoint-baseline sensitivities are reported in Appendix~\ref{app:fidelity}; these comparisons describe aggregate annotations rather than mechanisms. Restricting to the 129 items where baseline, A2, and the zero-start vector all reach 2,048 tokens gives cosine 0.067 [$-0.410$, $0.571$]. This post-treatment selection is not a causal duration adjustment, but shows substantial sensitivity to the observation window and population. Similar directions do not establish equivalent achieved profiles: magnitudes and individual coordinates differ.

At 8,192 tokens, natural-completion accuracy is 43.7\% unsteered, 73.6\% with A2, 67.3\% with zero-start RL, and 81.9\% with A2-initialized RL. At 32,768, the corresponding values are 86.4\%, 93.5\%, 97.2\%, and 94.5\%. Mean token use falls under all three interventions; these operating points do not establish a budget-independent best arm.

The historical forced-answer probe gives the saved zero-start RL vector a 7.3 pp [3.3, 11.6] advantage at a 1,024-token reasoning cutoff, where A2 does not show a comparable gain. This probe supplies a cue and additional decoding; it retains an entire final answer once thinking has closed, even if that answer extends beyond the cutoff. It therefore compares answer-elicitation protocols, not accuracy at identical total computation, and does not resolve the completion explanation for the domain-vector matrix. Appendix~\ref{app:elicitation} reports the completed common-total-budget control, including matched SEAL and concise prompting.


